% Sections 4.1 to 4.4, from lectures/L04/appendix/a_theory.md.

\section{Vad är ett ekvationssystem?}
\label{c4:sec:ekvationssystem}
Ett \textbf{linjärt ekvationssystem} består av flera ekvationer med flera obekanta som ska
uppfyllas \emph{samtidigt}. Det vanligaste är ett system med \textbf{två ekvationer och två
obekanta}:
\[
  \begin{cases}
    a_1 x + b_1 y = c_1 \\
    a_2 x + b_2 y = c_2
  \end{cases}
\]
\textbf{Lösningen} är det par $(x, y)$ som uppfyller båda ekvationerna. Grafiskt är det
\textbf{skärningspunkten} mellan de två räta linjerna.

\subsection{Antal lösningar}
\label{c4:sec:antal}

\begin{narrowtable}{@{}ll@{}}
  \thead{Situation} & \thead{Antal lösningar} \\
  \midrule
  Linjerna skär varandra & Exakt en lösning \\
  Linjerna är parallella (men ej sammanfallande) & Ingen lösning \\
  Linjerna sammanfaller & Oändligt många lösningar \\
\end{narrowtable}

\section{Substitutionsmetoden}
\label{c4:sec:substitution}
Substitutionsmetoden går till så här:
\begin{enumerate}
  \item Välj en ekvation och uttryck en obekant med hjälp av den andra.
  \item Substituera (sätt in) uttrycket i den andra ekvationen.
  \item Lös den resulterande ekvationen, som har en obekant.
  \item Beräkna den andra obekanta.
  \item Kontrollräkna i \emph{båda} de ursprungliga ekvationerna.
\end{enumerate}

\begin{exempel}
Lös systemet
\[
  \begin{cases}
    x + 2y = 7 \\
    3x - y = 7
  \end{cases}
\]
\textbf{Steg 1:} Ur ekvation (1) fås $x = 7 - 2y$.
\par\noindent\textbf{Steg 2:} Sätt in i ekvation (2):
\[
  3(7 - 2y) - y = 7
\]
\textbf{Steg 3:} Förenkla och lös:
\[
  21 - 6y - y = 7 \quad \Rightarrow \quad 21 - 7y = 7 \quad \Rightarrow \quad y = 2
\]
\textbf{Steg 4:} Beräkna $x$:
\[
  x = 7 - 2 \cdot 2 = 3
\]
\textbf{Steg 5:} Kontroll: $3 + 4 = 7$~\checkmark{} och $9 - 2 = 7$~\checkmark
\par\noindent Lösningen är $(x, y) = (3, 2)$.
\end{exempel}

\section{Additionsmetoden (eliminationsmetoden)}
\label{c4:sec:addition}
\textbf{Idé:} Multiplicera ekvationerna med lämpliga tal så att en obekant får \emph{lika stora men
motsatta} koefficienter, och addera sedan ekvationerna för att eliminera den obekanten.
\begin{enumerate}
  \item Multiplicera en eller båda ekvationerna med lämpliga konstanter.
  \item Addera ekvationerna rad för rad, så att en obekant försvinner.
  \item Lös den kvarstående ekvationen.
  \item Beräkna den eliminerade obekanten med en av de ursprungliga ekvationerna.
\end{enumerate}

\begin{exempel}
Lös systemet
\[
  \begin{cases}
    2x + 3y = 12 \\
    4x - y = 10
  \end{cases}
\]
\textbf{Steg 1:} Multiplicera ekvation (2) med $3$:
\[
  \begin{cases}
    2x + 3y = 12 \\
    12x - 3y = 30
  \end{cases}
\]
\textbf{Steg 2:} Addera ekvationerna:
\[
  14x = 42 \quad \Rightarrow \quad x = 3
\]
\textbf{Steg 3:} Sätt in $x = 3$ i ekvation (1):
\[
  6 + 3y = 12 \quad \Rightarrow \quad y = 2
\]
Lösningen är $(x, y) = (3, 2)$.
\end{exempel}

\section{Tillämpning: Kirchhoffs lagar}
\label{c4:sec:kirchhoff}
Ekvationssystem är grundläggande för kretsteori. Med Kirchhoffs lagar kan ström- och
spänningsfördelningen i en krets beräknas.

\subsection{Kirchhoffs spänningslag (KVL)}
\label{c4:sec:kvl}
Summan av alla spänningar runt en sluten slinga är noll:
\[
  \sum U = 0
\]

\subsection{Kirchhoffs strömlag (KCL)}
\label{c4:sec:kcl}
Summan av strömmarna in i en nod är noll, det vill säga strömmarna in är lika med strömmarna ut:
\[
  \sum I = 0
\]

\begin{exempel}
Betrakta en krets med två slingor och strömmarna $I_1$ och $I_2$. KVL ger
\[
  \begin{cases}
    10 = 2I_1 + 4(I_1 - I_2) \\
    0 = 4(I_2 - I_1) + 6I_2
  \end{cases}
\]
vilket förenklas till
\[
  \begin{cases}
    6I_1 - 4I_2 = 10 \\
    -4I_1 + 10I_2 = 0
  \end{cases}
\]
Ur ekvation (2) fås $I_1 = 2{,}5 I_2$. Insatt i ekvation (1) ger det
\begin{gather*}
  6 \cdot 2{,}5 I_2 - 4I_2 = 10 \quad \Rightarrow \quad I_2 = \frac{10}{11} \approx 0{,}91\,\text{A},
  \\
  I_1 = 2{,}5 \cdot \frac{10}{11} = \frac{25}{11} \approx 2{,}27\,\text{A}.
\end{gather*}
\end{exempel}
